\begin{theorem}\label{C:thirteen-all-stars}
No three-object category whose cross arrows form a star and
whose total number of arrows is at most thirteen has classifying
space homotopy equivalent to $\CP^2$.
\end{theorem}
\begin{proof}
Theorem~\ref{C:CP2-lower13} handles smaller totals.  Write
$m,n,p$ for the endomorphism orders and $s,t$ for the two nonempty
leg sizes, so $m+n+p+s+t=13$.  The relative-product argument of
Lemma~\ref{C:twelve-all-stars} excludes the case in which all
three vertex spaces have zero rational $H^2$.

A vertex with nonzero rational $H^2$ has order at least five.
If two vertices have that property, every size is forced:
the endomorphism orders are $5,5,1$ and both legs are singletons.
If the trivial vertex is central, apply
Lemma~\ref{C:trivial-central-fork}.  Otherwise one singleton leg
has a rationally acyclic opposite leaf, and
Lemma~\ref{C:star-singleton-join} gives zero products.

Assume there is exactly one vertex with nonzero rational $H^2$.
Each other vertex has order at most five; by the complete
classification in Theorem~\ref{E:all-six-types}, its classifying
space is rationally acyclic.  Suppose first that the distinguished
vertex is the center.  A singleton leg is again excluded by
Lemma~\ref{C:star-singleton-join}.  Otherwise $s,t\geq2$, and
the leaf orders have sum at most four.  If both are at most two,
their rational algebras are semisimple.  The exact-coinvariant
Tor calculation in Lemma~\ref{C:twelve-all-stars} applies verbatim:
the center has order at most seven, and its nonzero fourth homology
would persist in arbitrarily high degrees, a contradiction.

The remaining central case has leaf orders $3,1$, central order
five and leg sizes $2,2$.  The central monoid is necessarily the
five-element rectangular-band sphere monoid $P_2$, so its
classifying space is $S^2$.  Each two-object collage has zero
rational positive-degree products by the two-point assertion of
Lemma~\ref{C:two-object-cup-pruning}.  Mayer--Vietoris for their
union along $BP_2$ gives an injection
\[
 H^4(BC;\mathbb Q)\longrightarrow
 H^4(BC_1;\mathbb Q)\oplus H^4(BC_2;\mathbb Q),
\]
because $H^3(BP_2;\mathbb Q)=0$.  Every degree-two product
restricts to zero in both terms, so it is zero in $BC$ as well.

Finally suppose the distinguished vertex is a leaf, of order
$m\geq5$, and let $n$ be the central order and $s$ the size of
this distinguished leg.  Collapsing the rationally acyclic
center splits the star into its two leg cofibers by a rational
cohomology ring isomorphism.  Each cofiber is rationally equivalent,
through such a ring isomorphism, to its original collage.  The
other collage has two acyclic vertices and zero products.
Thus a nonzero square must occur in the distinguished collage,
whose rational homology must be bounded.

For $s\leq2$, Lemma~\ref{C:two-object-cup-pruning} excludes that
collage.  For $s\geq3$, the other leaf and leg contribute at least
two arrows, hence $m+n\leq8$.  If $n\leq2$, the all-act bounded
criterion of Lemma~\ref{C:eight-bounded-cups} applies.  Otherwise
the only profile is
\[
 \text{center }3,\qquad \text{distinguished leaf }5,
 \qquad \text{other leaf }1,
 \qquad (s,t)=(3,1).
\]
The five-element monoid is again exactly $P_2$.

The last profile is completely enumerated by
\path{code/c_thirteen_star_residual.py}.  It uses all seven ordinary
isomorphism classes of order-three monoids from the independent
census.  Every right action on the three-point leg is enumerated
up to state conjugacy, every left $P_2$-action in every relative
state position is retained, and every commuting pair is tested.
The singleton leg has forced actions.  Taking the opposite covers
inward stars.  The result is exactly $357$ commuting bisets:
$107$ have nonzero mod-two $H_1$, $196$ have the wrong mod-two
$H_2$ dimension, and $54$ have a zero degree-two square with an
explicit normalized cochain primitive.  The program verifies all
$111741$ typed associativity equations of these complete tables.
Every table, action and certificate is retained in
\path{code/outputs/c_thirteen_star_residual.json}.
There is no unexcluded case.  This finishes all star profiles.
\end{proof}
